Nesta seção se aborda os objetivos do relatório e as razões de sua elaboração.
A paginação do relatório técnico e/ou científicos inicia-se a partir desta
página. A introdução é uma seção primária.

Para relatório técnico e/ou científicos, a ABNT NBR 10719 recomenda que
utilize espaçamento simples, mas pode utilizar o espaçamento 1,5~cm (a
Coordenação de cada Curso da UFAPE deve padronizar para todo documento) e para
as margens: anverso, esquerda e superior de 3~cm e direita e inferior de 2~cm;
verso, direita e superior de 3~cm e esquerda e inferior de 2~cm.

Segue abaixo as normas da ABNT que são indispensáveis para a elaboração de
relatório técnico e/ou científicos:

\begin{itemize}
  \item ABNT NBR 6023, Informação e documentação -- Referências -- Elaboração.
  \item ABNT NBR 6024, Informação e documentação -- Numeração progressiva das
        seções de um documento escrito -- Apresentação.
  \item ABNT NBR 6027, Informação e documentação -- Sumário -- Apresentação ABNT
        NBR 6028, Informação e documentação -- Resumo -- Apresentação.
  \item ABNT NBR 6034, Informação e documentação -- Índice -- Apresentação.
  \item IBGE. Normas de apresentação tabular. 3. ed. Rio de Janeiro: IBGE, 1993.
        Disponível em:
        \url{https://biblioteca.ibge.gov.br/visualizacao/livros/liv23907.pdf}.
        Acesso em: 1 fev. 2022.
\end{itemize}

Para a elaboração das seções primárias, secundárias, terciárias, quaternárias
e quinárias, a autor deve seguir as orientações da ABNT NBR 6024 que trata
sobre a numeração progressiva. Esta norma destaca gradualmente os títulos das
seções e orienta na utilização dos recursos de negrito, itálico ou sublinhado
e outros. As seções que estão no sumário devem estar idênticas às seções que
estão no texto (ABNT, 2012).

\textopadrao

\textopadrao
